
\section{Setup and notation}

Let $\cP \to U$ be a $G$-torsor over $U = C \setminus \mdls$, where $G$ is a
finite abelian group, $H \trianglelefteq G$, and $C$ is a smooth projective
curve. Fix $p_\infty \in \mdls$ and an integer $d$.
Assume the $G$ torsor $\cP \to U$ has ramification strictly bounded by $d$ at $p_\infty$. (i.e. $< d$)
according to the definition in \Cref{definition:local_ramification_boundness}
(i.e. the higher d'th-ramification group of the extension of DVRs at $p_\infty$ is trivial).

Recall the four quotient schemes, defined in:
\Cref{subsection:quotient_torsors_towers}
\[
\cP^{[d]_H},\qquad \cP^{[d]_G},\qquad (\cP_{G/H})^{(d)},\qquad (\cP_{G/H})^{[d]},
\]
together with the canonical Cartesian square
\[
\begin{tikzcd}[row sep=large, column sep=large]
\cP^{[d]_H} \ar[r,"H"] \ar[d]
  & (\cP_{G/H})^{(d)} \ar[d,"q"] \\
\cP^{[d]_G} \ar[r,"H"'] \ar[dr,"G"',bend right=15]
  & (\cP_{G/H})^{[d]} \ar[d,"G/H"] \\
  & U^{(d)},
\end{tikzcd}
\]
in which the arrows labelled $H$, $G/H$ and $G$ are torsors under the
indicated finite constant groups (in particular, \emph{\'etale}), and the
upper square is Cartesian (\Cref{prop:torsor_tower_diagram}).

The $G/H$-torsor $\cP_{G/H} \to U$ compactifies to a finite map of smooth
projective curves $f \colon C' \to C$. Pick $p'_\infty \in f^{-1}(p_\infty)$

\[
f^{(d)} \colon C'^{(d)} \to C^{(d)}
\]
for the induced map of $d$-fold symmetric powers; it sends
$d \cdot p'_\infty \mapsto d \cdot p_\infty$. Let
\[
\pi_C \colon X_C \to C^{(d)},\qquad \pi_{C'} \colon X_{C'} \to C'^{(d)}
\]
be the blowups at $d \cdot p_\infty$ and $d \cdot p'_\infty$ respectively, with
exceptional divisors $E_C,E_{C'}$ and generic points $\eta_C,\eta_{C'}$. The
DVRs of these generic points are
\[
V_C := \mathcal{O}_{X_C,\eta_C} \subset K(C^{(d)}),\qquad
V_{C'} := \mathcal{O}_{X_{C'},\eta_{C'}} \subset K(C'^{(d)}).
\]

\section{The master diagram}\label{sec:master}

The torsor tower over the open base $U^{(d)}$ embeds into the projective
geometry of $C^{(d)}$ and $C'^{(d)}$, which in turn admit canonical blowups
at the points at infinity. We assemble all of this into a single commutative
diagram.

\[
\begin{tikzcd}[row sep=2.4em, column sep=2.2em]
\cP^{[d]_H} \ar[r,"H"] \ar[d]
  & (\cP_{G/H})^{(d)} \ar[d,"q"] \ar[r,hookrightarrow,"\iota'"]
  & C'^{(d)} \ar[dd,"f^{(d)}"]
  & X_{C'} \ar[l,"\pi_{C'}"'] \ar[dd,dashed,"f_X"] \\
\cP^{[d]_G} \ar[r,"H"'] \ar[dr,"G"',bend right=10]
  & (\cP_{G/H})^{[d]} \ar[d,"G/H"]
  & & \\
  & U^{(d)} \ar[r,hookrightarrow,"\iota"']
  & C^{(d)}
  & X_C \ar[l,"\pi_C"]
\end{tikzcd}
\]

Reading the diagram:
\begin{itemize}
  \item The \textbf{left block} (columns 1--2) is the torsor tower of
        \Cref{prop:torsor_tower_diagram}: $H$-torsors horizontally, the
        $G/H$-torsor on the lower right, and the composite $G$-torsor
        $\cP^{[d]_G} \to U^{(d)}$.
  \item The \textbf{middle block} (column 3) records the compactifications.
        The map $\iota$ is the canonical open immersion
        $U^{(d)} \hookrightarrow C^{(d)}$. The map $\iota'$ is the open
        immersion $(\cP_{G/H})^{(d)} \hookrightarrow C'^{(d)}$ obtained from
        compactifying $\cP_{G/H} \subset C'$. The vertical $f^{(d)}$ is the
        finite morphism of symmetric powers; restricting along $\iota'$
        recovers the canonical map $(\cP_{G/H})^{(d)} \to U^{(d)}$, which is
        the composite of $q$ with the $G/H$-torsor.
  \item The \textbf{right block} (column 4) records the blowups at the
        diagonal points at infinity. The dashed arrow
        $f_X \colon X_{C'} \dashrightarrow X_C$ is the rational lift of
        $f^{(d)}$, defined on the strict transform; on generic points of
        the exceptional divisors it sends $\eta_{C'} \mapsto \eta_C$ so we have
        induced extension of DVRs $V_C \hookrightarrow V_{C'}$         
\end{itemize}

\section{Unramifiedness of $\cP^{[d]_G}$ at $\eta_C$}

\begin{theorem}\label{thm:main}
Suppose
\begin{enumerate}
  \item[\textnormal{(H1)}] $(\cP_{G/H})^{[d]} \to U^{(d)}$ is unramified at
        $\eta_C$; and
  \item[\textnormal{(H2)}] $\cP^{[d]_H} \to (\cP_{G/H})^{(d)}$ is unramified
        at $\eta_{C'}$.
  \item [\textnormal{(H3)}]$G=Z/p^rZ$ is cyclic of prime power order.
\end{enumerate}
Then $\cP^{[d]_G} \to U^{(d)}$ is unramified at $\eta_C$.
\end{theorem}

\subsection{Proof Attempts}

\subsubsection{Attempt 1}
\begin{proof}[Proof attempt (Galois-closure version)]
\textit{Setup.}
Write
\[
K := K(U^{(d)}) = K(C^{(d)}),\quad
M := K\bigl((\cP_{G/H})^{[d]}\bigr),\quad
K' := K\bigl((\cP_{G/H})^{(d)}\bigr) = K(C'^{(d)}),
\]
\[
L := K(\cP^{[d]_G}),\qquad L_H := K(\cP^{[d]_H}).
\]
The Galois closure of all these fields over $K$ is $K(\cP^{\times d})$. Indeed,
$\cP^{[d]_G}, \cP^{[d]_H}, (\cP_{G/H})^{(d)}, (\cP_{G/H})^{[d]}$ are defined in
\Cref{subsection:quotient_torsors_towers} as the quotients of $\cP^{\times d}$
by the subgroups
\[
\KG \rtimes S_d,\quad K_H \rtimes S_d,\quad H^d \rtimes S_d,\quad
(H^d + \KG) \rtimes S_d
\]
of the big group $\Gamma := G^d \rtimes S_d$ acting on $\cP^{\times d}$; and
$U^{(d)}$ is the quotient by $G^d \rtimes S_d$. So
\[
N := K(\cP^{\times d}) / K
\quad\text{is Galois with group } \Gamma,
\]
and the various subfields correspond to the subgroups above
(\Cref{prop:torsor_tower_diagram}). All five intermediate fields above are
fixed by the respective subgroups; in particular, the bottom map
$U^{(d)}\!\leftarrow\!N$ has Galois group $\Gamma$, with normal subgroups giving
the Galois subextensions
\[
\Gal(L/K) = G^d \rtimes S_d / \KG \rtimes S_d \;\cong\; G,
\qquad
\Gal(M/K) \cong G/H,
\]
\[
\Gal(L_H/K') \cong H,\qquad \Gal(K'/K) \cong (G/H)^d \rtimes S_d / S_d
\;\text{(not a group quotient if $K'/K$ is non-Galois),}
\]
together with the canonical surjections
$\Gal(L/K)\twoheadrightarrow\Gal(M/K)$ (summation modulo $H$) and
$\Gal(K(\cP^{\times d})/K')\twoheadrightarrow\Gal(L_H/K')$ (summation).

\smallskip
\textit{Goal in inertia language.} Fix a discrete valuation $\widetilde{W}$ on
$N$ extending $V_C$, chosen so that $\widetilde{W}|_{K'} = V_{C'}$ — this is
possible because both $V_{C'}$ and any extension of $V_C$ to $N$ restrict to
$V_C$ on $K$, and the Galois group $\Gamma$ acts transitively on the extensions.
Let
\[
J := I(\widetilde{W}/V_C) \;\subseteq\; \Gamma = G^d \rtimes S_d
\]
be the inertia group of $\widetilde{W}$ over $V_C$. Writing elements of
$\Gamma$ as $((g_1,\dots,g_d),\sigma)$, the surjection
$\Gamma \to G$ inducing $\Gal(L/K) = G$ is the summation
$((g_i),\sigma) \mapsto \sum_i g_i$. We must show
\begin{equation}\label{eq:goal-inertia}
\sum_{i=1}^d g_i = 0 \quad\text{in } G,
\quad\text{for every } ((g_i),\sigma) \in J.
\end{equation}
Equivalently, $J \subseteq \KG \rtimes S_d$. (Then the residue extension at the
image of $\widetilde{W}$ in $L$ is also separable: this is automatic from the
analogous separability statements in (H1) and (H2), which we suppress.)

\smallskip
\textit{What the hypotheses give.}

(H1) says $M/K$ is unramified at the restriction $\widetilde{W}|_M$. Under
the surjection $\Gamma \to G/H$, $((g_i),\sigma)\mapsto \sum_i g_i \bmod H$,
inertia maps to inertia, so
\begin{equation}\label{eq:H1-inertia}
\sum_{i=1}^d g_i \in H \quad\text{for every } ((g_i),\sigma) \in J.
\end{equation}

(H2) says $L_H/K'$ is unramified at $V_{C'} = \widetilde{W}|_{K'}$. Now
$\Gal(N/K') = H^d \rtimes S_d$, and the surjection
$H^d \rtimes S_d \twoheadrightarrow H$ is again summation. The inertia of
$\widetilde{W}$ over $V_{C'}$ in $\Gal(N/K')$ equals
$J \cap (H^d \rtimes S_d)$. So (H2) reads
\begin{equation}\label{eq:H2-inertia}
J \cap \bigl(H^d \rtimes S_d\bigr) \;\subseteq\; K_H \rtimes S_d
\qquad\text{i.e. } \sum_i h_i = 0 \text{ whenever } ((h_i),\sigma)\in J,\; h_i\in H.
\end{equation}

Combining \eqref{eq:H1-inertia}, an element $((g_i),\sigma)\in J$ has
$\sum g_i \in H$, hence $(g_i) \in H^d + \KG$. So we can decompose
$g_i = h_i + k_i$ with $(h_i) \in H^d$ and $(k_i) \in \KG$ (the decomposition
is unique modulo $(k_0, -k_0)$ for $k_0 \in K_H$, but $\sum h_i$ is well-defined
modulo $K_H$, hence $\sum h_i$ as an element of $H$ is canonical). Goal
\eqref{eq:goal-inertia} becomes
\begin{equation}\label{eq:goal-restate}
\sum_i h_i = 0 \quad\text{for every } ((g_i),\sigma)\in J,
\quad\text{where } g_i = h_i + k_i \text{ as above}.
\end{equation}

\smallskip
\textit{The case $(k_i) = 0$ (purely $H$-part).} If $((g_i),\sigma)\in J$ with
$g_i \in H$, i.e.\ $j \in J \cap (H^d \rtimes S_d)$, then \eqref{eq:H2-inertia}
gives $\sum g_i = 0$, hence $\sum h_i = 0$. \checkmark

\smallskip
\textit{Where the argument stalls.} For a general $((g_i),\sigma)\in J$ with
nonzero $(k_i)\in\KG$, hypotheses (H1) and (H2) do \emph{not} suffice to force
$\sum h_i = 0$. Concretely, the failure mode is the following: a subgroup
\[
J \;\subseteq\; (H^d + \KG)\rtimes S_d
\]
satisfying \eqref{eq:H2-inertia} is precisely a subgroup all of whose elements
``purely in $H^d$'' lie in $K_H$. But $J$ can still contain
elements $((g_i),\sigma)$ where $g_i = h_i + k_i$ with $(k_i)\neq 0$ and
$\sum h_i \neq 0$ in $H$ — e.g.\ a ``diagonal'' subgroup
$\{(h, k(h)) : h \in H_0\} \subseteq H \oplus K_{G/H}$ projecting non-trivially
to both factors.

\textcolor{red}{\textbf{Missing input.} The hypotheses (H1) and (H2) alone are
insufficient. To close \eqref{eq:goal-restate} we need an additional input that
constrains the ``$K_{G/H}$-part'' of $J$, equivalently the inertia of the
non-Galois comparison map $q : (\cP_{G/H})^{(d)} \to (\cP_{G/H})^{[d]}$ at
$\widetilde{W}|_{K'} = V_{C'}$ over $\widetilde{W}|_M$. A sufficient
hypothesis would be (H3): the extension $K'/M$ is unramified at $V_{C'}$
over its restriction to $M$, i.e.\ $J \cap \Gal(N/M) \subseteq H^d \rtimes S_d$.
Combined with (H1) and (H2), (H3) forces $J\subseteq\KG\rtimes S_d$: any
$((g_i),\sigma)\in J$ has $g_i = h_i + k_i$ and (H3) gives $(k_i) = 0$, so
$(h_i)\in K_H$ by the case above, so $\sum g_i = 0$. However, (H3) is
\emph{not} automatic — \Cref{lemma:torsor_unramified_codim_one} establishes
unramifiedness of $q$ only at codimension-1 valuations of $C'^{(d)}$ supported
in the boundary, whereas $V_{C'}$ is the blowup-divisor valuation at the
maximally degenerate diagonal point $d\cdot p'_\infty\in C'^{(d)}$
(codimension $d$). In the tame cyclic test case $G/H = \Z/n$, a direct
computation in the formal model at $d\cdot p'_\infty$ shows that the inertia
of $V_{C'}$ in $K'/M$ has order $\gcd(d,n)$ — generally non-trivial — so
(H3) fails for tame compactifications when $\gcd(d,n)>1$. The truth of
\Cref{thm:main} in this regime therefore relies on a more subtle interaction
between the inertia structure at $V_C$ and the ramification bound on
$\cP\to U$ at $p_\infty$, not exploited above.}
\end{proof}
